\documentclass[12pt,letterpaper,oneside,onecolumn,openany]{book}
%% openany evita poner página adicional por capítulo
\usepackage[utf8]{inputenc}
\usepackage[spanish, es-tabla]{babel} %% es-tabla cambia cuadros por tablas
\usepackage{amsmath}
\usepackage{amsfonts}
\usepackage{amssymb}
%%%%%%%%%%%%%%%%%%%%%%%%%%%%%%%%%%%%
\usepackage{graphicx} %% gráficos
\DeclareGraphicsExtensions{.pdf,.jpeg,.png,.eps}
\usepackage{subfig}
\usepackage{float} %% PARA UBICAR FIGURAS
%%%%%%%%%TABLAS%%%%%%%%%%%%%%%%%%%%%
\usepackage{longtable}
\usepackage{booktabs} 
%\usepackage{tabls}
%\usepackage{ctable}
\usepackage{bigstrut}
\usepackage{multirow}
\usepackage{bigdelim}
\usepackage{rotating}
%\usepackage{rotfloat}
%%%% para páginas horizontales
\usepackage{pdflscape}
%%%%%%%%%%%%%%%%%%%%%%%%%%%%%%%%%%%%%%
\usepackage[sort&compress]{natbib} % para contraer referencias de texto
%%%%%%%%%%%%%%%%%%%%%%%%%%%%%%%%%%%%%%
\usepackage[pdftex,
            bookmarks,
            bookmarksnumbered=true,
            bookmarksopen=false,
            colorlinks,
            citecolor=blue,
            filecolor=blue,
            linkcolor=blue,
            urlcolor=blue,
            breaklinks=true, %% para títulos largos
            ]{hyperref}
\usepackage[left=4cm,right=2cm,top=3cm,bottom=2cm]{geometry}
\sloppy 
\pagestyle{plain} %% un solo número por página
\renewcommand{\baselinestretch}{1.0} %% interlineado de 1.0
\parskip=3mm %% espacio entre párrafos
\parindent=0mm %% espacio a principio de párrafo
%%http://minisconlatex.blogspot.com.co/2011/04/como-escribir-una-tesis-con-latex.html
%%%%%%%%%%% configuración de nomenclatura%%%%%%%%%%%%%%%%%%%%
\usepackage[intoc]{nomencl} % Nomenclature package
\makenomenclature
\usepackage{ifthen}
\renewcommand{\nomname}{Nomenclatura}
 \RequirePackage{ifthen}
 \renewcommand{\nomgroup}[1]
 {%
 	\ifthenelse{\equal{#1}{C}}{\item[\underline{Conjuntos:}]}
 	{%
 		\ifthenelse{\equal{#1}{P}}{\item[\underline{Parámetros:}]}
 		{
 			\ifthenelse{\equal{#1}{V}}{\item[\underline{Variables:}]}{}
 		}
 	}
 }
 %%%%%%%%%%%%%%%%%%%%%%%%%%%%%%%%%%%%%%%%%%%%%%%%%%%%%%%%%%%%%%
% \addto{\captionsspanish}{\def\chaptername{}}

\usepackage{titlesec}
\titleformat{\chapter}
	{\fontsize{17.28pt}{0pt}\bfseries}{\thechapter}{1em}{}
\titlespacing{\chapter}{1cm}{0pt}{30pt}
\titleformat{\section}
  {\normalfont\fontsize{14}{0pt}\bfseries}{\thesection}{1em}{}
\titlespacing{\section}{0pt}{24pt}{6pt}
\titleformat{\subsection}
  {\normalfont\fontsize{12}{0pt}}{\thesubsection}{1em}{}
\titlespacing{\subsection}{0pt}{24pt}{6pt}
%%%%%%%%%%%%%%%%%%%%%%GRAFICOS%%%%%%%%%%%%%%%%%%%%%%%%%%%%%%%%%
\usepackage{pgfgantt}
\usepackage{tikz}
%%%%%%%%%%%%%%%%%%%%%%%%%%%%%%%%%%%%%%%%%%%%%%%%%%%%%%%%%%%%%%%
\usepackage{fancyhdr}

\fancypagestyle{landscape}{
    \fancyhf{} % Limpia todos los encabezados y pies de página predefinidos
    \fancyfoot[C]{\thepage} % Número de página centrado en el pie de página
    \renewcommand{\headrulewidth}{0pt} % Elimina la línea del encabezado
    \renewcommand{\footrulewidth}{0pt} % Elimina la línea del pie de página
    \fancyfoot[C]{\makebox[\textwidth][c]{\begin{sideways}\thepage\end{sideways}}} % Colocar el número de página girado en el centro del pie de página
}
%%%%%%%%%%%%%%%%%%%%%%%%%%%%%%%%%%%%%%%%%%%%%%%%%%%%%%%%%%%%%%%
\renewcommand{\familydefault}{\rmdefault}			
%%%%%%%%%%%%%%%%%%%%%%%%%%%%%%%%%%%%%%%%%%%%%%%%%%%%%%%%%%%%%%
\begin{document}
%%%%%%%%%%%%%%%%%%%%%%%%%%%%%%%%%%%%%%%%%%%%%%%%%%%%%%%%%%%%%%
%%%%%%%%%%%%%%%%%%%%%%%%%%%%%%%%%%%%%%%%%%%%%%%%%%%%%%%%%%%%
% General data
%%%%%%%%%%%%%%%%%%%%%%%%%%%%%%%%%%%%%%%%%%%%%%%%%%%%%%%%%%%%

\def\ThesisTitle{NegALS: NEXT GENERATION ASSISTANT FOR THE EFFECTIVE MANAGEMENT OF LINUX SYSTEMS WITH PACKAGE SUGGESTIONS}

\def\Program{MASTER'S DEGREE IN SYSTEMS AND COMPUTER ENGINEERING}

\def\Year{2025}
\def\Date{March 27, 2025}

%%%%%%%%%%%%%%%%%%%%%%%%%%%%%%%%%%%%%%%%%%%%%%%%%%%%%%%%%%%%
% People
%%%%%%%%%%%%%%%%%%%%%%%%%%%%%%%%%%%%%%%%%%%%%%%%%%%%%%%%%%%%

\def\AuthorName{ENG. ANDR\'ES MANUEL PRIETO \'ALVAREZ}
\def\AdvisorName{MsC. LINA ELIZABETH PORRAS SANTANA}
\def\DirectorName{PhD. JOS\'E A. JARAMILLO VILLEGAS}

%%%%%%%%%%%%%%%%%%%%%%%%%%%%%%%%%%%%%%%%%%%%%%%%%%%%%%%%%%%%
% Stack
%%%%%%%%%%%%%%%%%%%%%%%%%%%%%%%%%%%%%%%%%%%%%%%%%%%%%%%%%%%%

\def\GenModel{llama3.2}
%%%%%%%%%%%%%%%%%%%%%%%%%%%%%%%%%%%%%%%%%%%%%%%%%%%%%%%%%%%%%%

%% Páginas no númeradas, portadas, dedicatoria, agredecimientos, introducción
\pagenumbering{gobble}
\begin{titlepage}
	\begin{center}
		%  \vspace{3cm}
		\hrule height2.5pt
		\vspace{.1cm}
		\hrule height1pt
		\vspace{20 pt}
		\begin{figure}[H]
			\begin{center}
				\includegraphics[width=5.15cm]{Figures/utp.pdf}
			\end{center}
		\end{figure}
		\large{\Program\\}
		\vspace{50 pt}
		\large{MASTER'S THESIS: \\}
		\large{\ThesisTitle\\}
		\vspace{50 pt}
		\large{AUTHOR: \\}
		\large{\AuthorName\\}


		\vfill
		\large{PEREIRA-\Year}
		\vspace{7 pt}
		\hrule height1pt
		\vspace{.1cm}
		\hrule height2.5pt
	\end{center}
\end{titlepage}

\newpage
\begin{titlepage}
	\thispagestyle{empty}
	\begin{center}

		\large{\ThesisTitle\\}
		\vspace{3cm}
		\large{\AuthorName\\}

		\vspace{3cm}
		\large{Degree project presented as a requirement to opt for the Master's degree in Systems and Computing Engineering\\}
		\vspace{3cm}
		\large{Director: \DirectorName\\}
		\large{Adviser: \AdvisorName\\}
		% \vspace{3cm}
		\vfill
		\large{UNIVERSIDAD TECNOLÓGICA DE PEREIRA\\}
		\large{\Program\\}
		\large{PEREIRA\\}
		\large{\Year}
	\end{center}
\end{titlepage}
\newpage
\begin{center}
	Year 2025
\end{center}
\begin{flushright}
	\begin{tabular}{r}
		\multicolumn{1}{c}{\textbf{ACCEPTANCE NOTE}} \\
		\\
		\textbf{} \bigstrut[b]                       \\
		\hline
		\textbf{} \bigstrut                          \\
		\hline
		\textbf{} \bigstrut                          \\
		\hline
		\bigstrut[t]                                 \\
		\\
		\\
		\\
		\\
		\textbf{} \bigstrut[b]                       \\
		\hline
		Director \bigstrut[t]                        \\
		\\
		\\
		\\
		\\
		\textbf{} \bigstrut[b]                       \\
		\hline
		Jury \bigstrut[t]                            \\
		\\
		\\
		\\
		\\
		\\
		\textbf{} \bigstrut[b]                       \\
		\hline
		Jury \bigstrut[t]                            \\
	\end{tabular}%

\end{flushright}
\vspace{96pt}
\begin{center}
	Pereira (\Date)
\end{center}

\chapter*{}
\vspace{1cm}
\begin{center}
	\textit{DEDICATION}
\end{center}
\vspace{1.5cm}
\begin{flushright}
	\textit{I dedicate this work to the free software community and the developer communities whose collaboration and generosity have enriched my professional and personal development. This work is a testimony of my gratitude to you.}
\end{flushright}
\vspace{1.0cm}
\begin{flushright}
	Eng. Andrés Manuel Prieto Álvarez.
\end{flushright}
\chapter*{}
\vspace{1cm}
\begin{center}
	\textit{ACKNOWLEDGMENTS}
\end{center}
\vspace{1.5cm}

\textit{I deeply thank my family for their unconditional support throughout all these years. Your meals, love, understanding and encouragement have been my greatest strength on this academic journey.}

\textit{I want to express my sincere gratitude to my thesis director José and my co-director Lina, for their invaluable guidance, support and dedication throughout this research process. Their experience and advice were fundamental for the development of this work.}

\textit{Finally, I would like to acknowledge all those who contributed in some way or another to this project. Specially Eng. Castañeda for helping in a rush as soon as the idea for "Lyra" came into the table and BA. Díaz for her endless corrections about my writing and grammar. Their encouragement and assistance were invaluable on this path towards the completion of this thesis.}
\vspace{1.0cm}
\begin{flushright}
	Eng. Andrés Manuel Prieto Álvarez.
\end{flushright}
\chapter*{Abstract}

GNU/Linux has become the operating system of choice for many users, especially in the software development and system management fields. However, its increasing complexity presents significant challenges \cite{gonzalez2005analyzing}. Traditionally designed for terminal-based interactions, Linux requires innovative solutions to bridge the gap between its powerful capabilities and accessibility for a broader audience. This research explores the feasibility of employing a Large Language Model \cite{sheng2023s-loraLLM, devlin2019bertLLM}, specifically \GenModel, to create a cost-effective Autonomous Agent that enhances the functionalities of a Virtual Assistant \cite{liu2023whenLLM, marukatat2021textLLM, hu2021loraLLM, sheng2023s-loraLLM, devlin2019bertLLM} tailored for GNU/Linux users. The goal is to improve user experience by providing intelligent, context-aware recommendations and assisting with software management.

Beyond usability concerns, GNU/Linux also faces structural challenges such as distribution fragmentation and hardware compatibility \cite{zambrano2022implementacion}, which can hinder adoption and performance consistency across different environments. Despite these challenges, its open-source nature and extensive toolset make it a preferred platform for AI-driven applications. The integration of Artificial Intelligence (AI) into GNU/Linux has enabled advancements in automation, predictive analytics, and system optimization, reinforcing its role in cutting-edge technological development. As automation continues to shape software management, leveraging AI to streamline system interactions becomes increasingly necessary.

In this context, this research investigates the potential of \GenModel\ to develop a Virtual Assistant specifically designed for GNU/Linux users firstly apt-based distro. By incorporating AI-driven adaptability, the assistant aims to bridge the gap between complex system operations and user-friendly interactions. Unlike conventional assistants, this model will not only address general queries, but also provide system-specific guidance tailored to the user's installed software, ultimately enhancing accessibility, efficiency, and overall user satisfaction.


\textbf{Keywords: GNU/Linux, Large language models (LLM), Autonomous Agent, Virtual Assistant} % Presentation
\pagenumbering{roman} %% númeración en romano
\renewcommand{\contentsname}{Table of Contents}
\tableofcontents
\renewcommand{\listfigurename}{Index of Figures}
\listoffigures
\renewcommand{\listtablename}{Index of Tables}
\listoftables
%%% cuerpo principal
\mainmatter %% cuerpo principal (la númeración es en arábigo)
\chapter*{Glossary}
Throughout this study, technical and specialized terms are used that may not be familiar to all readers. In order to facilitate understanding of the text and improve accessibility for a wide audience, this glossary has been included:
\begin{itemize}
	\item \textbf{GNU/Linux:} Is an operating system, a set of programs that allow you to interact with your computer and run other programs.
	\item \textbf{Large Language Model or "LLM":} Large language models (LLMs) are very large deep learning models that are pretrained on large amounts of data.
	\item \textbf{Virtual Assistant:} A virtual assistant is a software agent that helps users of computer systems, automating and performing tasks with minimal human-machine interaction.
	\item \textbf{SourceForge:} Is a collaboration website for software projects.
	\item \textbf{Wiki:} A wiki is a web-based collaborative platform that enables users to store, create and modify content in an organized manner.
	\item \textbf{Man Pages:} A man page (Short for "manual page") is a form of software documentation usually found on a Unix or Unix-like operating system.
\end{itemize}

%\include{Appex/Nomenclature}
\begin{titlepage}
	\begin{center}
		%  \vspace{3cm}
		\hrule height2.5pt
		\vspace{.1cm}
		\hrule height1pt
		\vspace{20 pt}
		\begin{figure}[H]
			\begin{center}
				\includegraphics[width=5.15cm]{Figures/utp.pdf}
			\end{center}
		\end{figure}
		\large{\Program\\}
		\vspace{50 pt}
		\large{MASTER'S THESIS: \\}
		\large{\ThesisTitle\\}
		\vspace{50 pt}
		\large{AUTHOR: \\}
		\large{\AuthorName\\}


		\vfill
		\large{PEREIRA-\Year}
		\vspace{7 pt}
		\hrule height1pt
		\vspace{.1cm}
		\hrule height2.5pt
	\end{center}
\end{titlepage}

\newpage
\begin{titlepage}
	\thispagestyle{empty}
	\begin{center}

		\large{\ThesisTitle\\}
		\vspace{3cm}
		\large{\AuthorName\\}

		\vspace{3cm}
		\large{Degree project presented as a requirement to opt for the Master's degree in Systems and Computing Engineering\\}
		\vspace{3cm}
		\large{Director: \DirectorName\\}
		\large{Adviser: \AdvisorName\\}
		% \vspace{3cm}
		\vfill
		\large{UNIVERSIDAD TECNOLÓGICA DE PEREIRA\\}
		\large{\Program\\}
		\large{PEREIRA\\}
		\large{\Year}
	\end{center}
\end{titlepage}
\newpage
\begin{center}
	Year 2025
\end{center}
\begin{flushright}
	\begin{tabular}{r}
		\multicolumn{1}{c}{\textbf{ACCEPTANCE NOTE}} \\
		\\
		\textbf{} \bigstrut[b]                       \\
		\hline
		\textbf{} \bigstrut                          \\
		\hline
		\textbf{} \bigstrut                          \\
		\hline
		\bigstrut[t]                                 \\
		\\
		\\
		\\
		\\
		\textbf{} \bigstrut[b]                       \\
		\hline
		Director \bigstrut[t]                        \\
		\\
		\\
		\\
		\\
		\textbf{} \bigstrut[b]                       \\
		\hline
		Jury \bigstrut[t]                            \\
		\\
		\\
		\\
		\\
		\\
		\textbf{} \bigstrut[b]                       \\
		\hline
		Jury \bigstrut[t]                            \\
	\end{tabular}%

\end{flushright}
\vspace{96pt}
\begin{center}
	Pereira (\Date)
\end{center}

\chapter*{}
\vspace{1cm}
\begin{center}
	\textit{DEDICATION}
\end{center}
\vspace{1.5cm}
\begin{flushright}
	\textit{I dedicate this work to the free software community and the developer communities whose collaboration and generosity have enriched my professional and personal development. This work is a testimony of my gratitude to you.}
\end{flushright}
\vspace{1.0cm}
\begin{flushright}
	Eng. Andrés Manuel Prieto Álvarez.
\end{flushright}
\chapter*{}
\vspace{1cm}
\begin{center}
	\textit{ACKNOWLEDGMENTS}
\end{center}
\vspace{1.5cm}

\textit{I deeply thank my family for their unconditional support throughout all these years. Your meals, love, understanding and encouragement have been my greatest strength on this academic journey.}

\textit{I want to express my sincere gratitude to my thesis director José and my co-director Lina, for their invaluable guidance, support and dedication throughout this research process. Their experience and advice were fundamental for the development of this work.}

\textit{Finally, I would like to acknowledge all those who contributed in some way or another to this project. Specially Eng. Castañeda for helping in a rush as soon as the idea for "Lyra" came into the table and BA. Díaz for her endless corrections about my writing and grammar. Their encouragement and assistance were invaluable on this path towards the completion of this thesis.}
\vspace{1.0cm}
\begin{flushright}
	Eng. Andrés Manuel Prieto Álvarez.
\end{flushright}
\chapter*{Abstract}

GNU/Linux has become the operating system of choice for many users, especially in the software development and system management fields. However, its increasing complexity presents significant challenges \cite{gonzalez2005analyzing}. Traditionally designed for terminal-based interactions, Linux requires innovative solutions to bridge the gap between its powerful capabilities and accessibility for a broader audience. This research explores the feasibility of employing a Large Language Model \cite{sheng2023s-loraLLM, devlin2019bertLLM}, specifically \GenModel, to create a cost-effective Autonomous Agent that enhances the functionalities of a Virtual Assistant \cite{liu2023whenLLM, marukatat2021textLLM, hu2021loraLLM, sheng2023s-loraLLM, devlin2019bertLLM} tailored for GNU/Linux users. The goal is to improve user experience by providing intelligent, context-aware recommendations and assisting with software management.

Beyond usability concerns, GNU/Linux also faces structural challenges such as distribution fragmentation and hardware compatibility \cite{zambrano2022implementacion}, which can hinder adoption and performance consistency across different environments. Despite these challenges, its open-source nature and extensive toolset make it a preferred platform for AI-driven applications. The integration of Artificial Intelligence (AI) into GNU/Linux has enabled advancements in automation, predictive analytics, and system optimization, reinforcing its role in cutting-edge technological development. As automation continues to shape software management, leveraging AI to streamline system interactions becomes increasingly necessary.

In this context, this research investigates the potential of \GenModel\ to develop a Virtual Assistant specifically designed for GNU/Linux users firstly apt-based distro. By incorporating AI-driven adaptability, the assistant aims to bridge the gap between complex system operations and user-friendly interactions. Unlike conventional assistants, this model will not only address general queries, but also provide system-specific guidance tailored to the user's installed software, ultimately enhancing accessibility, efficiency, and overall user satisfaction.


\textbf{Keywords: GNU/Linux, Large language models (LLM), Autonomous Agent, Virtual Assistant} % Introduction
\begin{titlepage}
	\begin{center}
		%  \vspace{3cm}
		\hrule height2.5pt
		\vspace{.1cm}
		\hrule height1pt
		\vspace{20 pt}
		\begin{figure}[H]
			\begin{center}
				\includegraphics[width=5.15cm]{Figures/utp.pdf}
			\end{center}
		\end{figure}
		\large{\Program\\}
		\vspace{50 pt}
		\large{MASTER'S THESIS: \\}
		\large{\ThesisTitle\\}
		\vspace{50 pt}
		\large{AUTHOR: \\}
		\large{\AuthorName\\}


		\vfill
		\large{PEREIRA-\Year}
		\vspace{7 pt}
		\hrule height1pt
		\vspace{.1cm}
		\hrule height2.5pt
	\end{center}
\end{titlepage}

\newpage
\begin{titlepage}
	\thispagestyle{empty}
	\begin{center}

		\large{\ThesisTitle\\}
		\vspace{3cm}
		\large{\AuthorName\\}

		\vspace{3cm}
		\large{Degree project presented as a requirement to opt for the Master's degree in Systems and Computing Engineering\\}
		\vspace{3cm}
		\large{Director: \DirectorName\\}
		\large{Adviser: \AdvisorName\\}
		% \vspace{3cm}
		\vfill
		\large{UNIVERSIDAD TECNOLÓGICA DE PEREIRA\\}
		\large{\Program\\}
		\large{PEREIRA\\}
		\large{\Year}
	\end{center}
\end{titlepage}
\newpage
\begin{center}
	Year 2025
\end{center}
\begin{flushright}
	\begin{tabular}{r}
		\multicolumn{1}{c}{\textbf{ACCEPTANCE NOTE}} \\
		\\
		\textbf{} \bigstrut[b]                       \\
		\hline
		\textbf{} \bigstrut                          \\
		\hline
		\textbf{} \bigstrut                          \\
		\hline
		\bigstrut[t]                                 \\
		\\
		\\
		\\
		\\
		\textbf{} \bigstrut[b]                       \\
		\hline
		Director \bigstrut[t]                        \\
		\\
		\\
		\\
		\\
		\textbf{} \bigstrut[b]                       \\
		\hline
		Jury \bigstrut[t]                            \\
		\\
		\\
		\\
		\\
		\\
		\textbf{} \bigstrut[b]                       \\
		\hline
		Jury \bigstrut[t]                            \\
	\end{tabular}%

\end{flushright}
\vspace{96pt}
\begin{center}
	Pereira (\Date)
\end{center}

\chapter*{}
\vspace{1cm}
\begin{center}
	\textit{DEDICATION}
\end{center}
\vspace{1.5cm}
\begin{flushright}
	\textit{I dedicate this work to the free software community and the developer communities whose collaboration and generosity have enriched my professional and personal development. This work is a testimony of my gratitude to you.}
\end{flushright}
\vspace{1.0cm}
\begin{flushright}
	Eng. Andrés Manuel Prieto Álvarez.
\end{flushright}
\chapter*{}
\vspace{1cm}
\begin{center}
	\textit{ACKNOWLEDGMENTS}
\end{center}
\vspace{1.5cm}

\textit{I deeply thank my family for their unconditional support throughout all these years. Your meals, love, understanding and encouragement have been my greatest strength on this academic journey.}

\textit{I want to express my sincere gratitude to my thesis director José and my co-director Lina, for their invaluable guidance, support and dedication throughout this research process. Their experience and advice were fundamental for the development of this work.}

\textit{Finally, I would like to acknowledge all those who contributed in some way or another to this project. Specially Eng. Castañeda for helping in a rush as soon as the idea for "Lyra" came into the table and BA. Díaz for her endless corrections about my writing and grammar. Their encouragement and assistance were invaluable on this path towards the completion of this thesis.}
\vspace{1.0cm}
\begin{flushright}
	Eng. Andrés Manuel Prieto Álvarez.
\end{flushright}
\chapter*{Abstract}

GNU/Linux has become the operating system of choice for many users, especially in the software development and system management fields. However, its increasing complexity presents significant challenges \cite{gonzalez2005analyzing}. Traditionally designed for terminal-based interactions, Linux requires innovative solutions to bridge the gap between its powerful capabilities and accessibility for a broader audience. This research explores the feasibility of employing a Large Language Model \cite{sheng2023s-loraLLM, devlin2019bertLLM}, specifically \GenModel, to create a cost-effective Autonomous Agent that enhances the functionalities of a Virtual Assistant \cite{liu2023whenLLM, marukatat2021textLLM, hu2021loraLLM, sheng2023s-loraLLM, devlin2019bertLLM} tailored for GNU/Linux users. The goal is to improve user experience by providing intelligent, context-aware recommendations and assisting with software management.

Beyond usability concerns, GNU/Linux also faces structural challenges such as distribution fragmentation and hardware compatibility \cite{zambrano2022implementacion}, which can hinder adoption and performance consistency across different environments. Despite these challenges, its open-source nature and extensive toolset make it a preferred platform for AI-driven applications. The integration of Artificial Intelligence (AI) into GNU/Linux has enabled advancements in automation, predictive analytics, and system optimization, reinforcing its role in cutting-edge technological development. As automation continues to shape software management, leveraging AI to streamline system interactions becomes increasingly necessary.

In this context, this research investigates the potential of \GenModel\ to develop a Virtual Assistant specifically designed for GNU/Linux users firstly apt-based distro. By incorporating AI-driven adaptability, the assistant aims to bridge the gap between complex system operations and user-friendly interactions. Unlike conventional assistants, this model will not only address general queries, but also provide system-specific guidance tailored to the user's installed software, ultimately enhancing accessibility, efficiency, and overall user satisfaction.


\textbf{Keywords: GNU/Linux, Large language models (LLM), Autonomous Agent, Virtual Assistant} % State of the Art
\begin{titlepage}
	\begin{center}
		%  \vspace{3cm}
		\hrule height2.5pt
		\vspace{.1cm}
		\hrule height1pt
		\vspace{20 pt}
		\begin{figure}[H]
			\begin{center}
				\includegraphics[width=5.15cm]{Figures/utp.pdf}
			\end{center}
		\end{figure}
		\large{\Program\\}
		\vspace{50 pt}
		\large{MASTER'S THESIS: \\}
		\large{\ThesisTitle\\}
		\vspace{50 pt}
		\large{AUTHOR: \\}
		\large{\AuthorName\\}


		\vfill
		\large{PEREIRA-\Year}
		\vspace{7 pt}
		\hrule height1pt
		\vspace{.1cm}
		\hrule height2.5pt
	\end{center}
\end{titlepage}

\newpage
\begin{titlepage}
	\thispagestyle{empty}
	\begin{center}

		\large{\ThesisTitle\\}
		\vspace{3cm}
		\large{\AuthorName\\}

		\vspace{3cm}
		\large{Degree project presented as a requirement to opt for the Master's degree in Systems and Computing Engineering\\}
		\vspace{3cm}
		\large{Director: \DirectorName\\}
		\large{Adviser: \AdvisorName\\}
		% \vspace{3cm}
		\vfill
		\large{UNIVERSIDAD TECNOLÓGICA DE PEREIRA\\}
		\large{\Program\\}
		\large{PEREIRA\\}
		\large{\Year}
	\end{center}
\end{titlepage}
\newpage
\begin{center}
	Year 2025
\end{center}
\begin{flushright}
	\begin{tabular}{r}
		\multicolumn{1}{c}{\textbf{ACCEPTANCE NOTE}} \\
		\\
		\textbf{} \bigstrut[b]                       \\
		\hline
		\textbf{} \bigstrut                          \\
		\hline
		\textbf{} \bigstrut                          \\
		\hline
		\bigstrut[t]                                 \\
		\\
		\\
		\\
		\\
		\textbf{} \bigstrut[b]                       \\
		\hline
		Director \bigstrut[t]                        \\
		\\
		\\
		\\
		\\
		\textbf{} \bigstrut[b]                       \\
		\hline
		Jury \bigstrut[t]                            \\
		\\
		\\
		\\
		\\
		\\
		\textbf{} \bigstrut[b]                       \\
		\hline
		Jury \bigstrut[t]                            \\
	\end{tabular}%

\end{flushright}
\vspace{96pt}
\begin{center}
	Pereira (\Date)
\end{center}

\chapter*{}
\vspace{1cm}
\begin{center}
	\textit{DEDICATION}
\end{center}
\vspace{1.5cm}
\begin{flushright}
	\textit{I dedicate this work to the free software community and the developer communities whose collaboration and generosity have enriched my professional and personal development. This work is a testimony of my gratitude to you.}
\end{flushright}
\vspace{1.0cm}
\begin{flushright}
	Eng. Andrés Manuel Prieto Álvarez.
\end{flushright}
\chapter*{}
\vspace{1cm}
\begin{center}
	\textit{ACKNOWLEDGMENTS}
\end{center}
\vspace{1.5cm}

\textit{I deeply thank my family for their unconditional support throughout all these years. Your meals, love, understanding and encouragement have been my greatest strength on this academic journey.}

\textit{I want to express my sincere gratitude to my thesis director José and my co-director Lina, for their invaluable guidance, support and dedication throughout this research process. Their experience and advice were fundamental for the development of this work.}

\textit{Finally, I would like to acknowledge all those who contributed in some way or another to this project. Specially Eng. Castañeda for helping in a rush as soon as the idea for "Lyra" came into the table and BA. Díaz for her endless corrections about my writing and grammar. Their encouragement and assistance were invaluable on this path towards the completion of this thesis.}
\vspace{1.0cm}
\begin{flushright}
	Eng. Andrés Manuel Prieto Álvarez.
\end{flushright}
\chapter*{Abstract}

GNU/Linux has become the operating system of choice for many users, especially in the software development and system management fields. However, its increasing complexity presents significant challenges \cite{gonzalez2005analyzing}. Traditionally designed for terminal-based interactions, Linux requires innovative solutions to bridge the gap between its powerful capabilities and accessibility for a broader audience. This research explores the feasibility of employing a Large Language Model \cite{sheng2023s-loraLLM, devlin2019bertLLM}, specifically \GenModel, to create a cost-effective Autonomous Agent that enhances the functionalities of a Virtual Assistant \cite{liu2023whenLLM, marukatat2021textLLM, hu2021loraLLM, sheng2023s-loraLLM, devlin2019bertLLM} tailored for GNU/Linux users. The goal is to improve user experience by providing intelligent, context-aware recommendations and assisting with software management.

Beyond usability concerns, GNU/Linux also faces structural challenges such as distribution fragmentation and hardware compatibility \cite{zambrano2022implementacion}, which can hinder adoption and performance consistency across different environments. Despite these challenges, its open-source nature and extensive toolset make it a preferred platform for AI-driven applications. The integration of Artificial Intelligence (AI) into GNU/Linux has enabled advancements in automation, predictive analytics, and system optimization, reinforcing its role in cutting-edge technological development. As automation continues to shape software management, leveraging AI to streamline system interactions becomes increasingly necessary.

In this context, this research investigates the potential of \GenModel\ to develop a Virtual Assistant specifically designed for GNU/Linux users firstly apt-based distro. By incorporating AI-driven adaptability, the assistant aims to bridge the gap between complex system operations and user-friendly interactions. Unlike conventional assistants, this model will not only address general queries, but also provide system-specific guidance tailored to the user's installed software, ultimately enhancing accessibility, efficiency, and overall user satisfaction.


\textbf{Keywords: GNU/Linux, Large language models (LLM), Autonomous Agent, Virtual Assistant} % Theorical Framework
\begin{titlepage}
	\begin{center}
		%  \vspace{3cm}
		\hrule height2.5pt
		\vspace{.1cm}
		\hrule height1pt
		\vspace{20 pt}
		\begin{figure}[H]
			\begin{center}
				\includegraphics[width=5.15cm]{Figures/utp.pdf}
			\end{center}
		\end{figure}
		\large{\Program\\}
		\vspace{50 pt}
		\large{MASTER'S THESIS: \\}
		\large{\ThesisTitle\\}
		\vspace{50 pt}
		\large{AUTHOR: \\}
		\large{\AuthorName\\}


		\vfill
		\large{PEREIRA-\Year}
		\vspace{7 pt}
		\hrule height1pt
		\vspace{.1cm}
		\hrule height2.5pt
	\end{center}
\end{titlepage}

\newpage
\begin{titlepage}
	\thispagestyle{empty}
	\begin{center}

		\large{\ThesisTitle\\}
		\vspace{3cm}
		\large{\AuthorName\\}

		\vspace{3cm}
		\large{Degree project presented as a requirement to opt for the Master's degree in Systems and Computing Engineering\\}
		\vspace{3cm}
		\large{Director: \DirectorName\\}
		\large{Adviser: \AdvisorName\\}
		% \vspace{3cm}
		\vfill
		\large{UNIVERSIDAD TECNOLÓGICA DE PEREIRA\\}
		\large{\Program\\}
		\large{PEREIRA\\}
		\large{\Year}
	\end{center}
\end{titlepage}
\newpage
\begin{center}
	Year 2025
\end{center}
\begin{flushright}
	\begin{tabular}{r}
		\multicolumn{1}{c}{\textbf{ACCEPTANCE NOTE}} \\
		\\
		\textbf{} \bigstrut[b]                       \\
		\hline
		\textbf{} \bigstrut                          \\
		\hline
		\textbf{} \bigstrut                          \\
		\hline
		\bigstrut[t]                                 \\
		\\
		\\
		\\
		\\
		\textbf{} \bigstrut[b]                       \\
		\hline
		Director \bigstrut[t]                        \\
		\\
		\\
		\\
		\\
		\textbf{} \bigstrut[b]                       \\
		\hline
		Jury \bigstrut[t]                            \\
		\\
		\\
		\\
		\\
		\\
		\textbf{} \bigstrut[b]                       \\
		\hline
		Jury \bigstrut[t]                            \\
	\end{tabular}%

\end{flushright}
\vspace{96pt}
\begin{center}
	Pereira (\Date)
\end{center}

\chapter*{}
\vspace{1cm}
\begin{center}
	\textit{DEDICATION}
\end{center}
\vspace{1.5cm}
\begin{flushright}
	\textit{I dedicate this work to the free software community and the developer communities whose collaboration and generosity have enriched my professional and personal development. This work is a testimony of my gratitude to you.}
\end{flushright}
\vspace{1.0cm}
\begin{flushright}
	Eng. Andrés Manuel Prieto Álvarez.
\end{flushright}
\chapter*{}
\vspace{1cm}
\begin{center}
	\textit{ACKNOWLEDGMENTS}
\end{center}
\vspace{1.5cm}

\textit{I deeply thank my family for their unconditional support throughout all these years. Your meals, love, understanding and encouragement have been my greatest strength on this academic journey.}

\textit{I want to express my sincere gratitude to my thesis director José and my co-director Lina, for their invaluable guidance, support and dedication throughout this research process. Their experience and advice were fundamental for the development of this work.}

\textit{Finally, I would like to acknowledge all those who contributed in some way or another to this project. Specially Eng. Castañeda for helping in a rush as soon as the idea for "Lyra" came into the table and BA. Díaz for her endless corrections about my writing and grammar. Their encouragement and assistance were invaluable on this path towards the completion of this thesis.}
\vspace{1.0cm}
\begin{flushright}
	Eng. Andrés Manuel Prieto Álvarez.
\end{flushright}
\chapter*{Abstract}

GNU/Linux has become the operating system of choice for many users, especially in the software development and system management fields. However, its increasing complexity presents significant challenges \cite{gonzalez2005analyzing}. Traditionally designed for terminal-based interactions, Linux requires innovative solutions to bridge the gap between its powerful capabilities and accessibility for a broader audience. This research explores the feasibility of employing a Large Language Model \cite{sheng2023s-loraLLM, devlin2019bertLLM}, specifically \GenModel, to create a cost-effective Autonomous Agent that enhances the functionalities of a Virtual Assistant \cite{liu2023whenLLM, marukatat2021textLLM, hu2021loraLLM, sheng2023s-loraLLM, devlin2019bertLLM} tailored for GNU/Linux users. The goal is to improve user experience by providing intelligent, context-aware recommendations and assisting with software management.

Beyond usability concerns, GNU/Linux also faces structural challenges such as distribution fragmentation and hardware compatibility \cite{zambrano2022implementacion}, which can hinder adoption and performance consistency across different environments. Despite these challenges, its open-source nature and extensive toolset make it a preferred platform for AI-driven applications. The integration of Artificial Intelligence (AI) into GNU/Linux has enabled advancements in automation, predictive analytics, and system optimization, reinforcing its role in cutting-edge technological development. As automation continues to shape software management, leveraging AI to streamline system interactions becomes increasingly necessary.

In this context, this research investigates the potential of \GenModel\ to develop a Virtual Assistant specifically designed for GNU/Linux users firstly apt-based distro. By incorporating AI-driven adaptability, the assistant aims to bridge the gap between complex system operations and user-friendly interactions. Unlike conventional assistants, this model will not only address general queries, but also provide system-specific guidance tailored to the user's installed software, ultimately enhancing accessibility, efficiency, and overall user satisfaction.


\textbf{Keywords: GNU/Linux, Large language models (LLM), Autonomous Agent, Virtual Assistant} % Result & discussion 
\begin{titlepage}
	\begin{center}
		%  \vspace{3cm}
		\hrule height2.5pt
		\vspace{.1cm}
		\hrule height1pt
		\vspace{20 pt}
		\begin{figure}[H]
			\begin{center}
				\includegraphics[width=5.15cm]{Figures/utp.pdf}
			\end{center}
		\end{figure}
		\large{\Program\\}
		\vspace{50 pt}
		\large{MASTER'S THESIS: \\}
		\large{\ThesisTitle\\}
		\vspace{50 pt}
		\large{AUTHOR: \\}
		\large{\AuthorName\\}


		\vfill
		\large{PEREIRA-\Year}
		\vspace{7 pt}
		\hrule height1pt
		\vspace{.1cm}
		\hrule height2.5pt
	\end{center}
\end{titlepage}

\newpage
\begin{titlepage}
	\thispagestyle{empty}
	\begin{center}

		\large{\ThesisTitle\\}
		\vspace{3cm}
		\large{\AuthorName\\}

		\vspace{3cm}
		\large{Degree project presented as a requirement to opt for the Master's degree in Systems and Computing Engineering\\}
		\vspace{3cm}
		\large{Director: \DirectorName\\}
		\large{Adviser: \AdvisorName\\}
		% \vspace{3cm}
		\vfill
		\large{UNIVERSIDAD TECNOLÓGICA DE PEREIRA\\}
		\large{\Program\\}
		\large{PEREIRA\\}
		\large{\Year}
	\end{center}
\end{titlepage}
\newpage
\begin{center}
	Year 2025
\end{center}
\begin{flushright}
	\begin{tabular}{r}
		\multicolumn{1}{c}{\textbf{ACCEPTANCE NOTE}} \\
		\\
		\textbf{} \bigstrut[b]                       \\
		\hline
		\textbf{} \bigstrut                          \\
		\hline
		\textbf{} \bigstrut                          \\
		\hline
		\bigstrut[t]                                 \\
		\\
		\\
		\\
		\\
		\textbf{} \bigstrut[b]                       \\
		\hline
		Director \bigstrut[t]                        \\
		\\
		\\
		\\
		\\
		\textbf{} \bigstrut[b]                       \\
		\hline
		Jury \bigstrut[t]                            \\
		\\
		\\
		\\
		\\
		\\
		\textbf{} \bigstrut[b]                       \\
		\hline
		Jury \bigstrut[t]                            \\
	\end{tabular}%

\end{flushright}
\vspace{96pt}
\begin{center}
	Pereira (\Date)
\end{center}

\chapter*{}
\vspace{1cm}
\begin{center}
	\textit{DEDICATION}
\end{center}
\vspace{1.5cm}
\begin{flushright}
	\textit{I dedicate this work to the free software community and the developer communities whose collaboration and generosity have enriched my professional and personal development. This work is a testimony of my gratitude to you.}
\end{flushright}
\vspace{1.0cm}
\begin{flushright}
	Eng. Andrés Manuel Prieto Álvarez.
\end{flushright}
\chapter*{}
\vspace{1cm}
\begin{center}
	\textit{ACKNOWLEDGMENTS}
\end{center}
\vspace{1.5cm}

\textit{I deeply thank my family for their unconditional support throughout all these years. Your meals, love, understanding and encouragement have been my greatest strength on this academic journey.}

\textit{I want to express my sincere gratitude to my thesis director José and my co-director Lina, for their invaluable guidance, support and dedication throughout this research process. Their experience and advice were fundamental for the development of this work.}

\textit{Finally, I would like to acknowledge all those who contributed in some way or another to this project. Specially Eng. Castañeda for helping in a rush as soon as the idea for "Lyra" came into the table and BA. Díaz for her endless corrections about my writing and grammar. Their encouragement and assistance were invaluable on this path towards the completion of this thesis.}
\vspace{1.0cm}
\begin{flushright}
	Eng. Andrés Manuel Prieto Álvarez.
\end{flushright}
\chapter*{Abstract}

GNU/Linux has become the operating system of choice for many users, especially in the software development and system management fields. However, its increasing complexity presents significant challenges \cite{gonzalez2005analyzing}. Traditionally designed for terminal-based interactions, Linux requires innovative solutions to bridge the gap between its powerful capabilities and accessibility for a broader audience. This research explores the feasibility of employing a Large Language Model \cite{sheng2023s-loraLLM, devlin2019bertLLM}, specifically \GenModel, to create a cost-effective Autonomous Agent that enhances the functionalities of a Virtual Assistant \cite{liu2023whenLLM, marukatat2021textLLM, hu2021loraLLM, sheng2023s-loraLLM, devlin2019bertLLM} tailored for GNU/Linux users. The goal is to improve user experience by providing intelligent, context-aware recommendations and assisting with software management.

Beyond usability concerns, GNU/Linux also faces structural challenges such as distribution fragmentation and hardware compatibility \cite{zambrano2022implementacion}, which can hinder adoption and performance consistency across different environments. Despite these challenges, its open-source nature and extensive toolset make it a preferred platform for AI-driven applications. The integration of Artificial Intelligence (AI) into GNU/Linux has enabled advancements in automation, predictive analytics, and system optimization, reinforcing its role in cutting-edge technological development. As automation continues to shape software management, leveraging AI to streamline system interactions becomes increasingly necessary.

In this context, this research investigates the potential of \GenModel\ to develop a Virtual Assistant specifically designed for GNU/Linux users firstly apt-based distro. By incorporating AI-driven adaptability, the assistant aims to bridge the gap between complex system operations and user-friendly interactions. Unlike conventional assistants, this model will not only address general queries, but also provide system-specific guidance tailored to the user's installed software, ultimately enhancing accessibility, efficiency, and overall user satisfaction.


\textbf{Keywords: GNU/Linux, Large language models (LLM), Autonomous Agent, Virtual Assistant} % Conclusions & Future work
\appendix
\mainmatter
\chapter{About results and Open-Source}
\label{apendice:lyra}

The results of this thesis are available at the Author's GitHub profile \url{https://github.com/AndresMpa} or the following repository \url{https://github.com/AndresMpa/lyra}; the results showed here should be replicable by following the instructions in the repository and replicating the hardware and software setup described in the Setup section at page \pageref{sec:setup}.

It must be mention that the author pretends to keep working on the repository, in order to replicate the results compiled here it could be necessary to search for the specific commit ID in which test \& measurements were taken.
\include{Appex/Appendix_b}
\bibliographystyle{apalike}
\bibliography{Bibliography}
\end{document}
%% Structure reference: https://repositorio.utp.edu.co/server/api/core/bitstreams/8ba3eb1f-2fae-462f-8c66-21e2aa11d2fc/content